% Efficient/Naive Lambda Calculus Encoding 
% Rewriting Guidance: None, sketch-guided, detailed rewriting 

\begin{figure}[b]
    \centering
\begin{tikzpicture}[scale=1.5]
    \tikzstyle{every node}=[font=\small\itshape, align=center]
    %
    \coordinate (origin)        at (0,0);
    \coordinate (top-left)      at (0,2.25);
    \coordinate (bottom-right)  at (5,0);
    \coordinate (top-right)     at (4.75,2);
    %
    \draw [->, thick, color={rgb,255:red,125;green,125;blue,125},line width=1pt] (origin) -- (top-right);
    \draw [->, thick] (origin) -- (top-left);
    \draw [->, thick] (origin) -- (bottom-right);
    %
    \node[anchor=north east] at (top-left) {sketch-guided\\ equality saturation};
    \node[anchor=east,yshift=-.5em] at ($(origin)!0.5!(top-left)$) {\normalfont{rewriting guidance}\\(\cref{sketching})};
    \node[anchor=south east] at (origin) {equality saturation};
    %
    \node[anchor=north west] at (origin) {naive};
    \node[anchor=north] at ($(origin)!0.5!(bottom-right)$) {$\lambda$~\normalfont{calculus encoding}\\(\cref{eqsat-bindings})};
    \node[anchor=north east] at (bottom-right) {efficient};
    %
    \node[fill=white] at ($(origin)!0.5!(top-right)$) {\bfseries scaling to complex optimizations\\\bfseries of functional programs};
\end{tikzpicture}
\caption{This paper is about scaling equality saturation to complex optimizations of functional programs by combining an efficient $\lambda$ calculus encoding with sketch-guided equality saturation.}
    \label{fig:overview}
\end{figure}

\section{Introduction}

Term rewriting has been effective in optimizing compilers for decades \cite{dershowitz1993-rewrite-systems}.
More recently, functional array languages like \Lift{}~\cite{lift-rewrite-2015} and \Rise{}~\cite{hagedorn2020-elevate} produce high-performance code for diverse hardware by using rewrite rules to define and explore optimization spaces.
However, deciding when to apply each rewrite rule, the so-called \emph{phase ordering problem}, is hard and has a huge impact on the performance of the rewritten program.
The challenge is that the global benefit of applying a rewrite rule depends on future rewrites.
Maximizing local benefit in a greedy fashion is not sufficient in the absence of a convergence property, i.e. confluence and termination, as local optima may be far away from global optima.
% \todo{and presumably most/many compiler rewrites don't have this convergence property?}

Rewriting strategies \cite{visser1998-strategies} allow programmers to control when to apply each rewrite rule, step-by-step.
Prior work on \Elevate{}~\cite{hagedorn2020-elevate} has shown that rewriting strategies can achieve \emph{complex optimizations} of \Rise{} programs in less than a second of compilation time.
However, these complex optimizations emerge from thousands of rewrite steps making the rewriting strategies challenging to write, as the phase ordering problem is passed on to the programmer.

% SPACE FOR AUTHOR NAMES
\vspace{8em}

Equality saturation \cite{tate2009-equality-saturation, willsey2021-egg} mitigates the phase ordering problem by exploring many possible ways to apply rewrite rules.
Starting from an input program, equality saturation grows an equality graph (\emph{e-graph}) by applying all possible rewrites iteratively until reaching a fixed point (saturation), achieving a performance goal, or timing out.
An e-graph efficiently represents a large set of equivalent programs, and is grown by repeatedly applying rewrite rules in a purely additive way.
Instead of replacing the matched left-hand side of a rewrite rule by its right-hand side, the equality between the left-hand side and the right-hand side is recorded in the e-graph.
After growing the e-graph, the best program found is extracted from it using a cost model, e.g. one that selects the fastest program.

Unfortunately, equality saturation does not scale to complex optimizations such as the ones applied to \Rise{} with \Elevate{}, producing huge e-graphs that exceed the memory available in most machines.
\textbf{To scale equality saturation to complex optimizations of functional programs}, this paper makes advances in two directions, as shown by the two axes in~\cref{fig:overview}.

\textbf{Rewriting Guidance.}
On each equality saturation iteration, the e-graph tends to grow larger as every possible rewrite rule is applied in a purely additive way. 
The growth rate is extremely rapid for some combinations of useful rewrite rules, like associativity and commutativity.
This makes exploring long rewrite sequences requiring many iterations unfeasible.
One way to address this issue is to limit the number of rules applied \cite{wang2020-spores, willsey2021-egg}, but this risks not finding optimizations that require an omitted rule.
An alternative is to use an external solver to speculatively add equivalences \cite{nandi2020-synthesizing-CAD}, but this requires the identification of sub-tasks that can benefit from being delegated.

This paper proposes \emph{sketch-guided equality saturation} that factors complex optimizations into a sequence of smaller optimizations, each sufficiently simple to be found by equality saturation.
The programmer guides rewriting by describing how a program
should evolve during optimization, through a sequence of \emph{sketches}: program patterns that leave details unspecified.
While sketches have previously been used as a starting point for program synthesis \cite{solar2008-synthesis-sketching},
our work uses sketches in a novel way as checkpoints to guide program optimization.

\textbf{Lambda Calculus Encoding.}
As almost all programming languages use variables, and hence name binding, this paper explores practical ways of efficiently encoding languages with name bindings for the purposes of equality saturation.
We focus on encoding the lambda calculus as it is the standard formalism for functional languages such as \Rise{}. Previously, equality saturation has used a naive lambda calculus encoding that is simple but inefficient \cite{willsey2021-egg}: the size of the e-graph quickly blows up.
The inefficiency can be avoided by explicitly avoid name bindings in the rewritten language~\cite{smith2021-access-patterns}. 
We show that lambda calculus name bindings can be managed by using De Bruijn indices to avoid overloading the e-graph with $\alpha$-equivalent terms, and that an approximate implementation of substitution avoids overloading the e-graph with intermediate substitution steps.

\smallskip

Our evaluation demonstrates that combining sketch-guided equality saturation with an efficient lambda calculus encoding enables complex optimizations of \Rise{} functional programs, that require tens of thousands of rewrite steps in term rewriting.
We first evaluate the effectiveness of our lambda calculus encoding techniques for three rewrite goals. Thereafter, we show that sketch guiding enables seven complex optimizations of matrix multiplication to be applied.
\Elevate{} rewriting strategies apply these complex  optimizations in less than a second with low memory consumption, but require ordering thousands of rewrite rules.
Unguided equality saturation abstracts over rewrite ordering, but can only locate the simplest of the seven optimizations before the e-graph exceeds the available memory.
Sketch-guided equality saturation applies all seven optimizations in seconds with low memory consumption, and only requires ordering up to three sketch guides.

\smallskip

To summarize, the main contributions of this paper are:
\begin{itemize}
  \item Proposing \emph{sketch-guided equality saturation}, a semi-automated process offering a novel, practical trade-off between rewriting strategies and equality saturation.
  The programmer guides multiple equality saturations by specifying a sequence of sketches describing how the program should evolve during optimization (\cref{sketching}).
  % Where rewriting strategies require fine-grained phase ordering (picking a sequence of rewrite rules), sketch-guided equality saturation only requires coarse-grained phase ordering (picking a sequence of sketches).
  % Where equality saturation does not scale due to an explosion of the e-graph size, sketch-guided equality saturation is able to scale further by dividing the optimization problem into smaller equality saturation searches.
  % as a new semi-automatic technique to perform complex optimizations that require long rewrite sequences, e.g. requiring more than \pwtcomment{how many rewrite steps? $10^{4}$?}, and cannot feasibly be discovered by equality saturation alone.
  
  % 
  \item Exploring new techniques to efficiently encode a polymorphically typed lambda calculus such as \Rise{} for the purpose of equality saturation.
  In particular, De Bruijn indices are used to avoid overloading the e-graph with $\alpha$-equivalent terms, and an approximate (\emph{extraction-based}) substitution is used to avoid overloading the e-graph with intermediate substitution steps (\cref{eqsat-bindings}).

 %  \item A discussion of how to support efficient equality saturation for the lambda calculus, as well as \kles{}, an concrete implementation of equality saturation for \Rise{}: a data-parallel functional language.
%  We demonstrate the efficiency of our implementation by optimizing a binomial filter application (\cref{eqsat-bindings}).
  % Practical choices to efficiently support equality saturation for the lambda calculus. These choices are realized in .
  % \kles{} broadens the applicability of equality saturation for languages with name bindings.
  % via DeBruijn indices and extraction based substitution 
  % This is a step towards scaling equality saturation to richer / more expressive languages.
  
  
  % \todo{a new application of program sketches to equality saturation to provide a semi-automatic technique for performing complex optimisations. That is, optimisations requiring long rewrite sequences that cannot be feasibly discovered by eq sat alone.}

  % Sketch-guided equality saturation broadens the applicability of equality saturation for long rewrite sequences.
  % offering trade-offs between control and automation of compiler optimizations
  % This is a step towards scaling equality saturation to more complex optimizations.
  % \item We present a novel beam extraction algorithm, which extracts the N best programs from an e-graph according to a given cost model, such that they satisfy a given program sketch.
  % , for multi-core processors 

  %\item An experimental evaluation demonstrating that sketch-guided equality saturation enables complex optimizations of a \Rise{} matrix multiplication requiring significantly less effort than using manually orchestrated rewriting.
  %While prior equality saturation is unable to apply the optimizations without running out of memory, sketch-guided equality saturation performs the optimization in seconds only requiring little manual guidance (\cref{evaluation}).

  \item Two systematic evaluations of \Rise{} applications.
  The effectiveness of our lambda calculus encoding is demonstrated by optimizing a binomial filter application (\cref{lambda-eval}).
  Sketch-guided equality saturation is evaluated by exploring seven complex optimizations of a matrix multiplication application, including loop blocking, vectorization, and multi-threading.
 These optimizations are infeasible with unguided equality saturation due to excessive runtime and memory consumption, but sketch-guided equality saturation performs the optimizations in seconds and with low memory consumption.
  No more than three sketch guides need to be written for any of the seven optimizations (\cref{mm-eval}).
  %are required to guide the search, i.e.\ significantly less effort than purely manual techniques. \todo{remove effort claim}
  %The techniques are realized in the \kles{} implementation for the \Rise{} data-parallel functional language.

  % \todo{
% \emph{A systematic comparative evaluation of sketch guided equality saturation for complex optimisations.} The exemplar is the optimisation of a RISE matrix multiplication using blocking, <name the optimisations>, and uses a sequence of YYY sketches. We demonstrate that SGES enables complex optimizations like <name the optimisations> that are not feasible with eq sat alone, and completes within XXXs. Moreover, a comparison with ELEVATE show that SGES requires significantly less effort than programmer orchestrated rewriting.<explicitly state, and quantify the effort metric here> (Section 5).
%   }
  % \todo{Consider introducing an acronym SGES for "sketch guided equality saturation"}
  % \todo{GC7 Currently the text of the abstract and introduction suggest that the only example you have is (a semi-realistic) matrix multiplication.

  % Suggest that you also state that you use other examples: a binomial filter, a harris (edge detector), ...}

  %Given a couple of sketches instead of a precise strategy, rewrite sequences are automatically discovered in seconds (\cref{evaluation}).
  % We leverage the existing \Rise{} language, a typed lambda calculus, and re-discover optimizations that were found in previous work by tediously orchestrating rewrite rules.
  % \todo{quantify}
\end{itemize}
